% !TeX program = xelatex
\documentclass[11pt,a4paper,UTF8,fontset=none]{ctexart}
\usepackage{easyR}

% 每次写作业，修改这里的信息和下面的正文即可。
\title{课程作业报告示例}
\author{姓名}
\reportsetup{course={课程名称},assignment={第一次作业},student-id={学号}}
% 如需日期，取消下一行注释。
% \date{\today}

\begin{document}
\maketitle

\section{如何使用这个模板}
本模板用于篇幅较短的中文课程报告。标题、姓名和学号在主文件中填写，
正文使用普通的章节、公式、图片和表格环境。编译时使用 XeLaTeX；
字体已随项目提供，不需要另行安装。

\subsection{公式与交叉引用}
例如，传输一个大小为 $S$ 字节的文件，在带宽为 $B$ bit/s 的网络中，
忽略协议开销时，理论传输时间为
\begin{equation}\label{eq:time}
  t=\frac{8S}{B}.
\end{equation}
通过 \verb|\label| 设置标签，再用 \verb|\autoref| 引用，
可以得到带名称的编号，如\autoref{eq:time}。文档增删内容后，编号会自动更新。

\subsection{表格}
表格建议使用三线表。\autoref{tab:example} 列出了两个计算结果。
\begin{table}[htbp]
  \centering
  \caption{1 MB 文件的理论传输时间}\label{tab:example}
  \begin{tabular}{@{}lcc@{}}
    \toprule
    文件大小 & 网络带宽 & 理论传输时间\\
    \midrule
    1 MB & 1 Mbps & 8 秒\\
    1 MB & 10 Mbps & 0.8 秒\\
    \bottomrule
  \end{tabular}
\end{table}

\subsection{图片}
图片放在 \texttt{figures/} 目录下，使用相对路径引用。
\autoref{fig:logo} 使用了仓库原有的校徽资源。
\begin{figure}[htbp]
  \centering
  \includegraphics[width=0.18\linewidth]{figures/sjtu_logo_notitle.pdf}
  \caption{插图示例}\label{fig:logo}
\end{figure}

\section{其他常用功能}
公式、定理、代码高亮、可分页文本框和子图均已默认启用，
直接使用相应环境即可。示例见 \texttt{examples/features.tex}。

\end{document}
